Mobile GUI agents sometimes must type text they saw many screens earlier. We call such a step a \emph{delayed information need}; it is absent from AITZ (0 of 18,643 steps) and present in 31\% of sampled GUI-Odyssey trajectories, with a median gap of 5 screens. Mining such a need requires a verifier: the typed text appears in the OCR (optical character recognition) of an earlier screen and not on the current one. We show that the same verifier supplies training labels for free, by marking every earlier screen as holding the text or not, and that this label-free supervision is what makes look-back work. Trained on these labels, a full-history reader localizes a source screen in 65.4\% of 188 verified cases against 54.8\% prompted, and a per-screen scorer reaches 72.3\% against 50.0\% prompted. The two trained variants differ by 6.9 points overall ($p=0.086$) and 2.1 end-to-end ($p=0.45$). The scorer's advantage concentrates at gaps of 8--12 screens (78.2\% vs.\ 56.4\%, $p=0.003$), where one call must weigh many candidates. On a fourth sample of trajectories held out from all design decisions, the scorer localizes 79.6\%, the trained full-history reader 76.3\% ($p=0.63$), and prompted full history 59.1\% ($p=0.003$); read with history raises end-to-end accuracy over prompted full history from 23.7\% to 49.5\% and 51.6\% ($p<0.001$ for both). We also report where both trained variants fail, and two evaluation filters that any study of this need requires.
